\documentclass[9pt,a4paper,twocolumn]{ctexart}

% --- Packages ---
\usepackage[top=1.2cm, bottom=1.2cm, left=1.2cm, right=1.2cm, columnsep=0.8cm]{geometry}
\usepackage{hyperref}
\usepackage{graphicx}
\usepackage{float}
\usepackage{longtable}
\usepackage{tabularx}
\usepackage{booktabs}
\usepackage{enumitem}
\usepackage{xcolor}
\usepackage{fancyhdr}
\usepackage{listings}
\usepackage{fontspec}
\usepackage{titlesec}
\usepackage{caption}

% --- Code block styling ---
\definecolor{codebg}{RGB}{245,245,245}
\definecolor{codeframe}{RGB}{220,220,220}
\definecolor{codekw}{RGB}{0,0,192}
\definecolor{codecomment}{RGB}{0,128,0}
\definecolor{codestring}{RGB}{163,21,21}
\definecolor{codenum}{RGB}{128,0,128}
\definecolor{codepunct}{RGB}{90,90,90}
\definecolor{badgegreen}{RGB}{34,139,34}
\definecolor{badgeblue}{RGB}{30,144,255}
\definecolor{badgeblack}{RGB}{50,50,50}

% --- Difficulty badges (rounded corners via tikz, pre-rendered in saveboxes) ---
\usepackage{tikz}
\newsavebox{\badgechubox}
\savebox{\badgechubox}{\tikz[baseline]{\node[fill=badgegreen,text=white,rounded corners=2.5pt,inner xsep=4pt,inner ysep=1.5pt,font=\sffamily\footnotesize\bfseries]{初级};}}
\newsavebox{\badgezhongbox}
\savebox{\badgezhongbox}{\tikz[baseline]{\node[fill=badgeblue,text=white,rounded corners=2.5pt,inner xsep=4pt,inner ysep=1.5pt,font=\sffamily\footnotesize\bfseries]{中级};}}
\newsavebox{\badgegaobox}
\savebox{\badgegaobox}{\tikz[baseline]{\node[fill=badgeblack,text=white,rounded corners=2.5pt,inner xsep=4pt,inner ysep=1.5pt,font=\sffamily\footnotesize\bfseries]{高级};}}
\newcommand{\lvchu}{\usebox{\badgechubox}}
\newcommand{\lvzhong}{\usebox{\badgezhongbox}}
\newcommand{\lvgao}{\usebox{\badgegaobox}}
\pdfstringdefDisableCommands{%
  \def\lvchu{[初级]}%
  \def\lvzhong{[中级]}%
  \def\lvgao{[高级]}%
}

\lstset{
  basicstyle=\ttfamily\scriptsize,
  keywordstyle=\color{codekw}\bfseries,
  commentstyle=\color{codecomment}\itshape,
  stringstyle=\color{codestring},
  backgroundcolor=\color{codebg},
  frame=single,
  rulecolor=\color{codeframe},
  breaklines=true,
  breakatwhitespace=false,
  breakindent=0pt,
  postbreak=\mbox{\textcolor{red}{$\hookrightarrow$}\space},
  numbers=left,
  numberstyle=\tiny\color{gray},
  columns=flexible,
  keepspaces=true,
  showstringspaces=false,
  tabsize=2,
  aboveskip=4pt,
  belowskip=4pt,
  extendedchars=true,
  literate={，}{{，}}1 {；}{{；}}1 {！}{{！}}1 {？}{{？}}1 {“}{{“}}1 {”}{{”}}1 {（}{{（}}1 {）}{{）}}1 {《}{{《}}1 {》}{{》}}1,
}

% --- Extra language definitions (no built-in listings support) ---
\lstdefinelanguage{json}{
  morekeywords={true,false,null},
  sensitive=true,
  morestring=[b]",
  comment=[l]{//},
  morecomment=[s]{/*}{*/},
}
\lstdefinelanguage{yaml}{
  morekeywords={true,false,null,yes,no,on,off},
  sensitive=false,
  morecomment=[l]{\#},
  morestring=[b]",
  morestring=[d]',
}
\lstdefinelanguage{http}{
  morekeywords={GET,POST,PUT,DELETE,PATCH,HEAD,OPTIONS,HTTP,Host,Accept,Authorization,Content-Type,Content-Length,User-Agent},
  sensitive=true,
  morecomment=[l]{\#},
  morestring=[b]",
}
\lstdefinelanguage{TypeScript}{
  morekeywords={abstract,any,as,async,await,break,case,catch,class,const,continue,debugger,declare,default,delete,do,else,enum,export,extends,false,finally,for,from,function,get,if,implements,import,in,instanceof,interface,keyof,let,module,namespace,never,new,null,number,of,package,private,protected,public,readonly,return,set,static,string,super,switch,symbol,this,throw,true,try,type,typeof,undefined,unknown,var,void,while,with,yield,Record,Array,Promise,AsyncIterable,ReadableStream,Date,Error,Object,JSON},
  sensitive=true,
  morecomment=[l]{//},
  morecomment=[s]{/*}{*/},
  morestring=[b]",
  morestring=[b]',
  morestring=[b]`{`},
}

% --- Compact spacing ---
\linespread{0.95}
\setlength{\parskip}{1pt plus 1pt minus 1pt}
\setlength{\parindent}{0pt}

% --- Table styling ---
\renewcommand{\arraystretch}{1.1}

% --- Heading styling (numbered) ---
\titleformat{\section}{\normalsize\bfseries}{\thesection\quad}{0em}{}
\titlespacing{\section}{0pt}{6pt}{2pt}
\titleformat{\subsection}{\small\bfseries}{\thesubsection\quad}{0em}{}
\titlespacing{\subsection}{0pt}{4pt}{1pt}
\titleformat{\subsubsection}{\small\bfseries}{\thesubsubsection\quad}{0em}{}
\titlespacing{\subsubsection}{0pt}{3pt}{1pt}

% --- Page style ---
\pagestyle{fancy}
\fancyhf{}
\fancyhead[L]{\small\emph{\leftmark}}
\fancyhead[R]{\small\thepage}
\renewcommand{\headrulewidth}{0.4pt}
\renewcommand{\sectionmark}[1]{}  % prevent sections from overwriting leftmark

% --- Hyperref setup ---
\hypersetup{
  colorlinks=true,
  linkcolor=blue,
  urlcolor=blue,
  citecolor=blue,
}

% --- Caption format ---
\DeclareCaptionFormat{myformat}{\small\bfseries #1#2#3}
\captionsetup{format=myformat}

\begin{document}

\title{万字详解 RAG 基础概念}
\date{}
\maketitle
\markboth{TypeScript 版 Agent 知识}{ TypeScript 版 Agent 知识}

\section*{什么是 RAG？}


\textbf{RAG（Retrieval-Augmented Generation，检索增强生成）} 就是把信息检索和大语言模型绑在一起用。系统先从知识库里检索出和当前问题相关的片段，知识库可以是数据库、文档集合，也可以是企业内部系统。然后把这些片段和原始问题一起喂给 LLM，让模型基于检索内容回答，而不是只靠训练时记住的知识。



\begin{figure}[H]
\centering
\includegraphics[width=0.85\textwidth]{/home/mao/demos/latex_demos/ai-agent-knowledge/HTML/万字详解 RAG 基础概念 _ JavaGuide_files/rag-simplified-architecture-diagram.jpg}
\caption*{RAG 示意图}
\end{figure}



\section*{为什么需要 RAG？}


\begin{figure}[H]
\centering
\includegraphics[width=0.85\textwidth]{/home/mao/demos/latex_demos/ai-agent-knowledge/HTML/万字详解 RAG 基础概念 _ JavaGuide_files/rag-llm-challenges.jpg}
\caption*{RAG（检索增强生成）如何解决 LLM 的核心挑战}
\end{figure}



LLM 训练数据再大，也绕不开几个问题。RAG 正好可以在这些地方进行弥补。



\textbf{第一是知识时效性。}



预训练模型的知识会停在训练数据截止时间点。训练后发生的新事件、新政策、新产品文档，模型默认是不知道的，除非通过联网、工具调用或外部知识注入来补。RAG 的做法是动态检索外部知识源，把最新的相关内容直接送给 LLM，让它不用只依赖参数里的旧知识。



\textbf{第二是私有数据访问。}



企业内部的产品文档、知识库、客户数据，不可能让公开 LLM 随便访问。RAG 在用户提问时只提取和问题相关的片段给 LLM，不需要暴露全部数据，模型也能基于企业自己的知识回答。



\textbf{第三是幻觉问题。}



LLM 会编造事实。RAG 通过提供明确参考文本，让模型尽量基于证据回答，能降低幻觉概率，但不能彻底消除幻觉。检索错误、上下文噪声、引用错配、模型不遵循指令，都可能导致错误答案。生产级 RAG 通常还要配引用校验、答案评估、拒答机制和人工反馈闭环。



\section*{RAG 的常见用途有哪些？}


RAG 最适合“答案依赖外部资料，并且资料会变化或很长”的场景。它先从知识库里检索相关内容，再让大模型基于检索结果生成回答，减少胡编，同时提高可追溯性。



常见场景包括这些：


\begin{itemize}[itemsep=1pt, topsep=2pt]
  \item 客服机器人：基于产品知识库做问答、排障、流程引导，比如“如何退换货”“某型号设备报错码怎么处理”。
  \item 研发 / 运维 Copilot：检索代码库、接口文档、告警手册，辅助定位问题和生成修复建议。
  \item 医疗助手：检索指南、药品说明、院内规范后生成辅助建议，但不做最终诊断，比如“某药禁忌是什么”“依据指南解释检查指标含义”。
  \item 法律咨询：基于法规条文、案例、合同模板检索，生成条款解释和风险提示。
  \item 教育辅导：从教材、讲义、题库中检索知识点，生成讲解和例题步骤。
  \item 企业内部助手：连接制度、SOP、会议纪要、技术文档，做检索、总结、对比。
  \item 投研、合规、审计、销售方案支持：处理报告、披露、内控、产品手册、标书模板等资料。
\end{itemize}


\section*{为什么有些企业还是宁愿用传统搜索而不是 RAG？}


不是所有问题都值得上 RAG。很多企业保留传统搜索，不是因为不知道 RAG 好用，而是用户需求本来就没到“生成答案”这一步。



如果用户只是想找一份制度原文、某个接口文档、一个合同模板，搜索框反而更直接。输入关键词，返回文档列表，用户自己点开确认，链路短、成本低、结果也更可控。RAG 则要先检索，再组织上下文，最后交给 LLM 生成答案。只要经过生成，就会多出延迟、Token 成本和总结偏差的风险。



所以选传统搜索还是 RAG，先看用户到底想要什么：是“帮我找到材料”，还是“帮我读完材料并给出结论”。


{\footnotesize
\begin{tabular}{|p{\dimexpr 0.150\columnwidth-2\tabcolsep\relax}|p{\dimexpr 0.350\columnwidth-2\tabcolsep\relax}|p{\dimexpr 0.500\columnwidth-2\tabcolsep\relax}|}
\hline
\textbf{维度} & \textbf{传统搜索（搜索框）} & \textbf{RAG（检索 + 生成）} \\
\hline
用户目标 & 找到文档、页面、附件 & 直接得到可读答案、总结或对比结论 \\
\hline
延迟与成本 & 极低，容易扩展 & 更高，需要检索和 LLM 推理 \\
\hline
可控性 / 可审计 & 强，直接给原文链接 & 弱一些，可能误解或总结偏差，需要引用与评测 \\
\hline
风险 & 低，主要是召回排序问题 & 更高，包括幻觉、引用错误、越权泄露 \\
\hline
数据治理 & 相对成熟，ACL、字段过滤都好做 & 更复杂，需要检索过滤、上下文脱敏、日志治理 \\
\hline
适用场景 & 编号、标题、关键词检索，找模板、找制度原文 & 客服解答、技术排障、制度解读、跨文档总结对比 \\
\hline
最佳实践 & ES / BM25 + 权限过滤 & 混合检索 + 重排 + 引用溯源 + 权限过滤 + 评测闭环 \\
\hline
\end{tabular}
}



实际落地时，很多企业会同时保留两套入口：\textbf{简单查找走搜索，复杂问答走 RAG}。这个组合通常比“所有问题都交给 RAG”更稳，也更省钱。



\section*{RAG 工作原理了解吗？}


RAG 的工程链路通常分两个阶段：离线索引和在线检索生成。索引阶段把原始文档处理成可检索的数据结构；在线阶段在用户提问时完成查询理解、检索召回、上下文构建和答案生成。



索引和检索阶段的简化流程图如下：



\begin{figure}[H]
\centering
\includegraphics[width=0.85\textwidth]{/home/mao/demos/latex_demos/ai-agent-knowledge/HTML/万字详解 RAG 基础概念 _ JavaGuide_files/rag-rag-engineering-link.jpg}
\caption*{索引和检索阶段的简化流程图}
\end{figure}



索引阶段主要做这些事：


\begin{enumerate}[itemsep=1pt, topsep=2pt]
  \item 输入文档：文本文件、PDF、网页、数据库记录都可以，只要有内容。
  \item 清理文档：去掉 HTML 标签、特殊字符等噪声。
  \item 增强文档：补充元数据，比如时间戳、分类标签，为后续检索提供过滤维度。
  \item 文档拆分（Chunking）：用文本分割器把文档切成较小片段。这一步要兼顾语义完整性、Embedding 模型输入长度、生成模型上下文窗口和召回粒度。Chunk 太大容易引入噪声，太小又可能丢上下文。拆分策略会直接影响召回质量，详细可以看 RAG 文档处理篇。
  \item 向量化表示（Embedding Generation）：通过嵌入模型将文本片段映射为语义向量，也就是高维稠密向量。常见嵌入模型包括 OpenAI 的 \texttt{text-embedding-3-small} / \texttt{text-embedding-3-large}，以及 Hugging Face 上的开源模型。
  \item 存储到向量存储或索引系统：把嵌入向量、原始内容和对应元数据存入向量存储或向量索引系统，比如 Milvus、pgvector、Elasticsearch / OpenSearch 向量检索，或基于 Faiss 构建本地向量索引。向量数据库选型、索引算法和 pgvector 实践可以看 RAG 向量库篇。
\end{enumerate}


索引过程通常离线完成。比如团队每周跑一次定时任务，把新增和变更的文档重新索引一遍。如果是用户上传文档这类动态场景，索引也可以在线完成，直接集成到主应用里。



检索是在线进行的。用户提问之后，系统通常会走下面这些步骤：


\begin{enumerate}[itemsep=1pt, topsep=2pt]
  \item 接收请求：拿到用户的自然语言查询。有些系统会先做查询改写或扩充，让后续检索更容易命中。
  \item 查询向量化：用嵌入模型把查询也转成向量，这样才能和文档向量在同一个空间里比较。
  \item 信息检索（R）：在向量库里做相似性搜索，把和查询向量最相关的文档片段捞出来。
  \item 上下文增强（A）：把检索片段、原始问题、系统指令和引用要求组织成 Prompt，交给 LLM。
  \item 输出生成（G）：LLM 输出自然语言回复，同时附上参考资料链接。
  \item 结果反馈（可选）：用户不满意时可以反馈，系统再调整 Prompt 或检索策略。有些实现也支持多轮对话来逐步完善回答。
\end{enumerate}


检索效果不稳定时，问题往往出在查询改写、召回策略、排序或上下文质量上。优化方向可以看 RAG 优化篇。



\section*{Embedding 是什么？}


Embedding 就是把文本变成一串数字。更准确地说，它会把文本映射到一个高维稠密向量空间里，让语义接近的文本在向量空间中距离更近。



比如这三句话：


\begin{itemize}[itemsep=1pt, topsep=2pt]
  \item “如何申请退款？”
  \item “退款流程是什么？”
  \item “订单怎么取消并退钱？”
\end{itemize}


它们字面不一样，但语义接近。好的 Embedding 模型会把它们映射到相近位置，向量检索才能把相关 Chunk 找出来。



\begin{figure}[H]
\centering
\includegraphics[width=0.85\textwidth]{/home/mao/demos/latex_demos/ai-agent-knowledge/HTML/万字详解 RAG 基础概念 _ JavaGuide_files/rag-2-embedding-map-text-to-semantic-space.jpg}
\caption*{Embedding：把文本映射到语义空间}
\end{figure}



Embedding 维度通常是 768、1024、1536、3072 等。维度越高，能表达的信息越丰富，但存储、索引和相似度计算成本也越高。以 OpenAI Embedding 为例，\texttt{text-embedding-3-small} 默认输出 1536 维，\texttt{text-embedding-3-large} 默认输出 3072 维，并支持通过 \texttt{dimensions} 参数降低输出维度。



常见 Embedding 模型可以分成两类：


{\footnotesize
\begin{tabular}{|p{\dimexpr 0.050\columnwidth-2\tabcolsep\relax}|p{\dimexpr 0.765\columnwidth-2\tabcolsep\relax}|p{\dimexpr 0.185\columnwidth-2\tabcolsep\relax}|}
\hline
\textbf{类型} & \textbf{代表模型} & \textbf{适合场景} \\
\hline
闭源 API & OpenAI \texttt{text-embedding-3-small} / \texttt{text-embedding-3-large}、Cohere Embed、Jina Embeddings API & 追求开箱即用、多语言效果、少运维 \\
\hline
开源模型 & BGE 系列、GTE 系列、E5 系列、Jina Embeddings 开源模型 & 数据不能出内网、需要私有化部署、希望控制成本 \\
\hline
\end{tabular}
}



选 Embedding 模型时，别只看榜单排名。MTEB（Massive Text Embedding Benchmark）可以作为参考，但最后还是要用自己的业务问题评测召回率、相关性和延迟。



Embedding 模型也不是“实时理解世界”的东西。它主要负责把文本映射到向量空间，能力重点是语义匹配。如果遇到非常新的术语、梗、产品名或领域缩写，仍然要通过业务语料评测确认召回效果。



\section*{向量相似度怎么计算？}


文本变成向量之后，检索系统还要判断哪个向量和查询最接近。常见相似度或距离度量有三种。


{\footnotesize
\begin{tabular}{|p{\dimexpr 0.403\columnwidth-2\tabcolsep\relax}|p{\dimexpr 0.169\columnwidth-2\tabcolsep\relax}|p{\dimexpr 0.429\columnwidth-2\tabcolsep\relax}|}
\hline
\textbf{度量方式} & \textbf{含义} & \textbf{特点} \\
\hline
余弦相似度（Cosine Similarity） & 看两个向量方向是否一致 & 对向量长度不敏感，RAG 场景最常用 \\
\hline
内积（Inner Product / Dot Product） & 看两个向量对应维度乘积之和 & 如果向量已经 L2 归一化，内积和余弦相似度在排序结果上通常等价 \\
\hline
欧氏距离（L2 Distance） & 看两个点在空间中的绝对距离 & 对向量幅度更敏感，适合模型或索引明确按 L2 训练 / 优化的场景 \\
\hline
\end{tabular}
}



RAG 常用余弦相似度的原因：RAG 关注的是语义方向是否接近，而不是向量长度本身；余弦相似度对长度不敏感，更适合文本语义检索。实际项目里还要和 Embedding 模型推荐的距离度量、向量库索引类型保持一致，否则可能导致索引无法命中或召回效果下降。



\section*{RAG 与传统搜索引擎的区别是什么？}


\begin{figure}[H]
\centering
\includegraphics[width=0.85\textwidth]{/home/mao/demos/latex_demos/ai-agent-knowledge/HTML/万字详解 RAG 基础概念 _ JavaGuide_files/rag-rag-vs-search-engine.jpg}
\caption*{RAG 与传统搜索引擎的区别}
\end{figure}



RAG 和传统搜索都在“找信息”，但拿到信息之后做的事不一样。



传统搜索拿到候选文档后，按相关性排好序，直接把结果列表给用户。每个结果彼此独立，用户自己点开、自己判断。它更像一个排序器。



RAG 会把检索到的多个知识片段一起放进 LLM 上下文，让模型做跨文档归纳和信息整合，最后生成一个直接能读的答案。它更像一个信息综合器。



几个差异比较关键：


\begin{enumerate}[itemsep=1pt, topsep=2pt]
  \item 检索机制：传统搜索主要靠倒排索引和关键词匹配，BM25 是经典算法；现代搜索系统也会加语义召回和重排。RAG 的检索方式更灵活，向量检索、BM25、混合检索、图检索、数据库查询都可以用，关键是检索结果要进入 LLM 上下文参与答案生成。
  \item 结果形态：搜索给文档列表，用户还要二次阅读；RAG 给答案，并尽量标出引用来源。
  \item 数据范围：传统搜索擅长全网爬虫和大规模索引；RAG 更常用于企业内部知识库和垂直领域，让 LLM 低成本获得特定领域知识补充。
  \item 成本和延迟：搜索响应快，成本可控；RAG 多了 LLM 推理，延迟和成本都会上去。
\end{enumerate}


\section*{RAG 和微调怎么选？}


RAG 与微调的选型问题很常见。



二者区分：RAG 解决模型缺少新知识或私有知识的问题，微调解决输出风格、格式和任务行为不符合要求的问题。



类比：员工手册经常需要查阅。RAG 是随查随用，把手册放在外部，每次回答前检索；微调是把手册内容内化进模型。手册频繁改版时，RAG 更新索引即可；微调需要重新准备数据、训练和评测，成本差异明显。


{\footnotesize
\begin{tabular}{|p{\dimexpr 0.048\columnwidth-2\tabcolsep\relax}|p{\dimexpr 0.337\columnwidth-2\tabcolsep\relax}|p{\dimexpr 0.614\columnwidth-2\tabcolsep\relax}|}
\hline
\textbf{维度} & \textbf{RAG} & \textbf{微调（Fine-tuning）} \\
\hline
知识更新 & 更新知识库或向量索引即可 & 通常需要重新准备数据并训练 \\
\hline
数据安全 & 知识保留在外部库，按需检索 & 训练样本中的模式和部分知识会固化到微调模型参数中，敏感数据进入训练流程前需要额外评估合规和数据治理要求 \\
\hline
幻觉控制 & 可引用原文，便于溯源和校验 & 模型仍可能编造，且引用来源不天然可见 \\
\hline
成本结构 & 检索成本 + 输入 Token 成本 + 向量库成本 & 数据标注、训练 GPU、评测和版本管理成本 \\
\hline
适合场景 & 知识密集型问答、企业知识库、法规制度、产品文档、实时信息 & 风格适配、格式控制、领域术语对齐、固定任务行为优化 \\
\hline
主要风险 & 检索不到、召回噪声、权限过滤复杂 & 数据过拟合、知识过期、训练和回滚成本高 \\
\hline
\end{tabular}
}



二者也可以结合。先用微调让模型更懂领域术语、输出格式和任务边界，再用 RAG 提供实时知识和可追溯证据。这类组合在客服、法律、医疗、金融投研等场景里很常见。



选型结论：知识变动频繁、需要引用来源，优先 RAG；输出风格和任务行为不稳定，考虑微调；既要懂领域表达又要查实时知识，可以两者结合。



不过这里有个现实限制：两者结合意味着两套系统都要维护，成本不低。团队资源有限时，先把 RAG 做稳，再考虑是否引入微调，通常更务实。



\section*{长上下文窗口会取代 RAG 吗？}


不会。



长上下文窗口确实让很多任务变简单了。比如把一整份报告丢进去，让模型从头读到尾，这类单文档深度分析很适合用长上下文。但它不等于可以把全部知识库都塞给模型。上下文越长，输入 Token 成本、首字延迟和推理噪声都会上升，效果未必更好。



长上下文适合的场景很明确：单篇长文档深度分析，一个代码仓库或一个项目目录的集中理解，长对话历史总结，或者一次性材料不多但需要完整阅读的任务。



知识库规模一大，长上下文就不够用了。企业知识库、客服工单、日志、合同库动辄百万到亿级文档片段，不可能每次都全塞进去。就算塞得进去，成本和延迟也扛不住。更麻烦的是，上下文里塞太多无关片段，模型反而更容易被噪声干扰，生成看起来完整但事实不稳的答案。“Lost in the Middle”问题说的就是这个，关键信息放在长上下文中间位置时更容易被忽略。



企业知识库还绕不开权限隔离。哪些内容用户能看，哪些不能看，不能靠“全塞进去”解决。RAG 可以在检索阶段做权限过滤，只把用户有权访问的内容放进上下文。长上下文做不了这件事。



还有一点经常被忽视：可追溯性。RAG 可以明确返回引用片段，审计时能溯源。长上下文把大量内容混在一起交给模型，用户很难判断回答到底基于哪段材料。



\section*{RAG 有哪些演进阶段？}


RAG 这两年一直在迭代，大致可以分成三个阶段。



\begin{figure}[H]
\centering
\includegraphics[width=0.85\textwidth]{/home/mao/demos/latex_demos/ai-agent-knowledge/HTML/万字详解 RAG 基础概念 _ JavaGuide_files/rag-2-evolution-stages.jpg}
\caption*{RAG 演进阶段}
\end{figure}


{\footnotesize
\begin{tabular}{|p{\dimexpr 0.133\columnwidth-2\tabcolsep\relax}|p{\dimexpr 0.589\columnwidth-2\tabcolsep\relax}|p{\dimexpr 0.278\columnwidth-2\tabcolsep\relax}|}
\hline
\textbf{阶段} & \textbf{典型链路} & \textbf{特点} \\
\hline
Naive RAG & 文档切块 → Embedding → Top-K 检索 → LLM 生成 & 最基础、最容易实现，适合 Demo 和简单知识库 \\
\hline
Advanced RAG & Query Rewrite / HyDE → 混合检索 → Rerank → 上下文压缩 → LLM 生成 & 重点解决召回不准、上下文噪声和排序不稳 \\
\hline
Modular RAG & 检索器、重排器、压缩器、路由器、生成器等模块可插拔组合 & 按业务场景动态路由，适合生产系统和复杂 Agent \\
\hline
\end{tabular}
}



Naive RAG 是起点，能跑通 Demo，但离生产通常还有距离。Advanced RAG 开始处理召回质量、噪声过滤和排序问题。Modular RAG 把各环节拆成可替换模块，更适合复杂场景。具体优化策略可以继续看 RAG 优化篇。



\section*{RAG 的核心优势和局限性是？}


先说优势。



\textbf{RAG 最大的好处是知识更新成本低。} 微调要重新准备数据、训练模型、评测效果，RAG 通常只需要更新知识库和索引。新闻、法规、产品文档这类经常变化的数据，用 RAG 维护起来会轻很多。



\textbf{它也能减少幻觉，并且方便追溯来源。} RAG 让模型从“凭记忆回答”变成“基于检索证据回答”。每个回答都可以挂到具体文档片段上，这在金融合规、医疗辅助、法律检索这些对准确性要求高的场景里很重要。当然，这不代表 RAG 就不会出错，检索错了、引用错了，答案一样会翻车。



\textbf{数据隔离也更容易做。} 检索层可实现多租户隔离和访问控制（ACL），确保用户只能看到权限范围内的数据。相比把敏感数据放进微调训练集，RAG 这套架构更适合做权限和合规治理。



\textbf{换领域的成本也低。} 不需要针对每个领域重新训练模型，把领域知识库建好、索引跑通，就能先用起来。



再看局限。RAG 不是银弹，坑也不少。



\textbf{检索质量决定上限。} GIGO 原则在这里特别明显：如果 Embedding 表达不准，或者分块策略把关键信息切丢了，召回内容和问题本身无关，下游 LLM 再强也救不回来。



\textbf{上下文也不是越长越好。} 虽然有些模型的 Context Window 已经扩展到百万级，但塞太多无关片段进去，模型注意力会被稀释，逻辑推理会被干扰，Token 开销也会跟着上升。



\textbf{延迟是另一个硬问题。} 完整链路要经过查询改写、向量化、相似度检索、重排序、上下文构建、LLM 生成，每一步都会增加耗时。对响应时间敏感的场景，不能只看答案质量，也要认真算延迟账。



\textbf{工程复杂度也不低。} 需要维护向量数据库，处理文档增量索引，持续优化检索策略，并完成权限过滤、引用溯源和评测闭环。相比直接调用 LLM API，RAG 的运维负担明显更重。



\textbf{Token 成本同样要算清楚。} RAG 省了训练成本，但每次请求都要带上下文，输入 Token 往往比普通对话高不少。文档片段塞得越多，账单和延迟都会一起涨。



\section*{⭐️ RAG 实战项目推荐}


参考项目：Spring Boot 4.0 + Java 21 + Spring AI + PostgreSQL + pgvector + RustFS + Redis，覆盖简历解析、岗位匹配、知识库 RAG。



项目地址：<https://github.com/Snailclimb/interview-guide> · <https://gitee.com/SnailClimb/interview-guide>



\section*{总结}


RAG 先从知识库检索相关内容，再让 LLM 基于检索结果回答。其价值不是提升模型能力，而是把回答拉回可检索、可引用、可审计的证据上。



关键点：


\begin{enumerate}[itemsep=1pt, topsep=2pt]
  \item RAG 主要解决的是 LLM 知识过时、碰不到私有数据、容易幻觉这几个问题。传统搜索给的是文档列表，RAG 给的是直接可读的答案；一个更像排序器，一个更像信息综合器。
  \item 知识变动频繁、需要引用来源时，优先考虑 RAG；如果要让模型按固定风格和格式输出，再考虑微调。
  \item 长上下文适合少量材料的深度分析，但企业级海量知识库、权限隔离和成本控制，还是要靠 RAG 这类检索链路来兜底。
\end{enumerate}


其局限同样明确：检索质量决定上限，上下文噪声会干扰生成，延迟、工程复杂度和 Token 成本真实存在。



Demo 跑通不等于生产可用；RAG 的难点通常不是接入向量库，而是持续评估和优化召回质量。


\end{document}
